\section{A sanctity of conversion and unfinished work}\label{sec:aug-synthesis}

Augustine's life has a strong narrative arc because he gave it one: restless
desire is gathered by grace into praise. Historical biography must both
receive and interrupt that arc. It receives the convert's testimony that God
acted through Scripture, friendship, sacrament, intellect, grief, and the
Church. It interrupts any temptation to make the narrator omniscient. His
partner's name is lost; opponents speak through unequal archives; coercive
policies affected persons outside the bishop's intended medicinal theory; and
the books later called timeless were made amid local emergencies.

The deepest continuity in the life is not a perfect doctrine possessed at
Milan. It is the practice of return: to Scripture after intellectual pride, to
common life after career, to the congregation after the study, to correction
after publication, and to penitential prayer at death. Augustine's sanctity,
as the Church receives it, includes extraordinary gifts of mind but is not
identical with brilliance. It is a converted use of those gifts in sacrament,
preaching, friendship, argument, service, and repentance.

His historical importance is correspondingly double. Few authors have opened
the interior life, biblical time, grace, ecclesial sacrament, and political
pilgrimage with comparable power. Few ancient bishops have supplied later
Western power with a more influential rationale for compelling religious
unity. A mature reception need not choose between gratitude and judgment.
Augustine himself teaches readers to revisit a corpus, name what needs
correction, and continue toward a truth no author owns.

Unresolved questions remain: the exact boundaries of his ancestry and language
world; details of the episcopal succession; the dating of some sermons and
letters; the textual history of the Rule; lost Donatist voices; the precise
relic itinerary; and how newly discovered texts may alter the balance again.
These limits do not empty the life. They keep a historical saint from becoming
an instrument made entirely in the reader's image.
